\scanpage{65}
\mspage{71}

\inserted{Il suffit de prouver que $f$ est une $\unsure{\mathbb D|\underline D}$-équiv.} \deleted{\unsure{Ceci s'applique}} \inserted{\unsure{Prenons}} p. ex. $f$ $\mathbb D$-propre. \deleted{Pour ceci,} \inserted{Il suffit} de prouver que si $Y$ est $\mathbb D$-contractile, alors $X$ l'est \deleted{que les $X_{/y_0}$ sont $\mathbb D$-contractiles. Or on} (quitte à l'appliquer \unsure{aux|au} $Y_{/y_0}$ \deleted{\unsure{au lieu de $Y$}}, après chgt de base \inserted{pour voir que les $X_{/y_0}$ sont contr.}).

\diagram{\begin{tikzcd}[ampersand replacement=\&]
\deleted{X_{/y_0}}\arrow[d,phantom,"{\deleted{\downarrow}}" description,"{\deleted{\unsure{f}}}"'] \& X\arrow[d,"f"']\\
\deleted{Y_{/y_0}} \& Y
\end{tikzcd}}

Pour ceci, il suffit de prouver que $f$ est une $\mathbb D$-équivalence, et on sait que ceci est le cas si pour tt \deleted{$F$} \uline{constant} sur

\deleted{\uline{Corollaire 3} Si $f$ est}

$Y$, \unsure{flèche|flèches} $F\xrightarrow{\mathrm{adj}} f_*f^*F$ est iso. En fait, c'est vrai pour \uline{tout} $F$ \unsure{;|:}

\uline{Corollaire 3} : Si $f:X\longrightarrow Y$ est \deleted{\unsure{co}} $\mathbb D$-propre ou $\mathbb D$-lisse, \inserted{à fibres $\mathbb D$-contractiles, alors} $f^*:\mathbb D(Y)\longrightarrow\mathbb D(X)$ est \uline{pl. fidèle}, i.e. les flèches d'adjonction

\[
 F\longrightarrow f_*f^*F,\qquad f_!f^*F\longrightarrow F
\]

dans $\mathbb D(Y)$ sont des iso.

Il suffit de le voir pour $f$ propre, et \unsure{cela|celà} résulte du calcul des fibres de $f_*f^*F$ \ldots

% slanted
\margin{\deleted{\unsure{aux seules fibres|aux fibres}}}

\deleted{\unsure{Prop.}} \uline{Prop. 5} : Les notions de foncteur $\mathbb D$-propre (resp. $\mathbb D$-lisse) est stable par changement de base (trivial), et par composition (utiliser Cor. 2). \unsure{Toute|Tout} \deleted{iso foncteur \illegible{1} fidèle est à la fois $\mathbb D$-propre et $\mathbb D$-lisse} \unsure{imm. fermée|sous-foncteur} est \deleted{\unsure{prop}} $\mathbb D$-propre, \unsure{tt} \unsure{imm. ouv.} $\mathbb D$-lisse. Plus gén, tout iso. local est $\mathbb D$-lisse; \deleted{tout} \unsure{co-}iso. local est $\mathbb D$-propre.

\uline{Prop. 6} Soient $X,Y$ sur $S$, $f:X\to Y$ un $S$-morphisme. Supposons $X$ et $Y$ propres sur $S$ \deleted{\unsure{lisses, et que}} \inserted{p.ex.} ou $X$ et $Y$ lisses sur $S$, et que $\forall s\in S$, $f_s:X_s\longrightarrow Y_s$ soit une $\mathbb D$-équiv, alors $f$ en est une $ {}^*$.

% margin note written sideways
\margin{$*$ \uline{\unsure{Dire|Disons}} : \deleted{Dans le cas $X,Y$ propres sur $S$} \unsure{au principal|au point que} $\equiv$ que sur la sous-catégorie des coeffts sur $Y$ qui proviennent de $\mathbb D(Y)$, $f^*:\mathbb D(Y)\longrightarrow\mathbb D(X)$ induit un foncteur pl. fidèle. Si $f$ propre, pr l'énoncé préc. \unsure{comme|car on a} $q_*q^*F\xrightarrow{\sim}p_*p^*F$ pour tt $F$, et cela résulte du calcul de \unsure{$p_*,q_*$|$q_*,q_*$} par passage aux fibres (\deleted{\unsure{$X,Y$}} \unsure{$q,p$} propres), et de l'hypothèse sur les $f_s$ OK.}
% archive sheet number: 64
